\documentclass{article}
\usepackage[utf8]{inputenc}
\usepackage{amsmath,amssymb}
\usepackage{graphicx}
\usepackage{hyperref}
\title{Synthesis of Biological and Architectural Forms via\\ GRA Nullification (Bio-Arch GRA)}
\author{O. Bits\footnote{Independent researcher, \texttt{https://orcid.org/0009-0004-1872-1153}}}
\date{\today}

\begin{document}
\maketitle

\begin{abstract}
We present an experiment in bio-inspired architectural design using the GRA (Geometric Recursive Analytics) multilevel nullification framework. A parametric column model with internal porosity is optimized to simultaneously satisfy biological criteria (lightweight, high specific strength) and architectural requirements (stability, buildability, aesthetic proportion). The optimization minimizes a sum of ``foam'' functionals derived from mechanical simulations and proportionality rules. Results demonstrate that GRA produces structurally efficient and visually harmonious forms resembling trabecular bone columns.
\end{abstract}

\section{Introduction}
GRA nullification \cite{gra_metatheorem} is a methodology for eliminating domain-specific artifacts by recursively canceling inter-domain discrepancies (foam). In previous work, it was applied to text-based synthesis (quantum poetry \cite{gra_quantum_poetry}). Here we extend the framework to geometric synthesis of biological and architectural forms.

\section{Methods}
\subsection{GRA Levels and Foam}
We define two levels:
\begin{align}
\Phi_{\text{bio}} &= w_1 (1 - F_{\text{cr}}/F_{\text{opt}})^2 + w_2 (m/m_{\text{opt}}-1)^2 + w_3 (p/p_{\text{opt}}-1)^2, \\
\Phi_{\text{arch}} &= v_1 (1 - F_{\text{cr}}/F_{\text{req}})^2 + v_2 \max(0,(t_{\text{min}}-t)/t_{\text{min}})^2 + v_3 (H/(2R)-\phi)^2,
\end{align}
where $F_{\text{cr}}$ is the Euler buckling load, $m$ mass, $p$ porosity, $t$ minimum wall thickness, and $\phi=1.618$. The total foam is $J = \lambda \Phi_{\text{bio}} + (1-\lambda)\Phi_{\text{arch}}$.

\subsection{Mechanical Simulation}
The column is a cylinder of radius $R$ and height $H$ with spherical pores. The critical load is $F_{\text{cr}} = \pi^2 E I_{\text{min}} / H^2$, where $I_{\text{min}}$ is the smallest area moment of inertia along the axis, computed analytically from the pore positions.

\subsection{Optimization}
We use a genetic algorithm (DEAP) with a genome encoding $R$, $H$, and up to 10 pores (active flag, coordinates, radius). Population size 50, 100 generations, crossover/mutation standard.

\section{Results}
After optimization, the best shape (Fig.~\ref{fig:column}) achieved $\Phi_{\text{bio}}=0.12$, $\Phi_{\text{arch}}=0.09$ for $\lambda=0.5$. The form exhibits a porosity of $0.47$, $F_{\text{cr}}$ within 90\% of target, and $H/D\approx 1.62$.

\begin{figure}[h]
\centering
\includegraphics[width=0.6\textwidth]{bio_column.png}
\caption{Optimized bio-arch column with internal pores.}
\label{fig:column}
\end{figure}

\section{Conclusion}
GRA nullification successfully merged biological and architectural design objectives, generating a lightweight, stable, and aesthetically pleasing column. Future work will replace the parametric model with a mesh variational autoencoder for richer topologies.

\bibliographystyle{plain}
\bibliography{references}
\end{document}
